\chapter{기본 도구와 세분의 경로 표현}
\label{chap:tools}

여기서부터 어려운 방향이다. \ref{chap:tools}--\ref{chap:main}장의 항목은 꼭짓점 집합 넓히기
(보조정리~\ref{lem:map_iff}), 작은 지지 집합의 평면성(보조정리~\ref{lem:small}),
연결성의 정의(정의~\ref{def:conn3}), 변 수축의 정의(정의~\ref{def:contract}),
순환 꼭짓점 수열의 정의(정의~\ref{def:cycseq}), 경로 체계의 정의(정의~\ref{def:paths})를 빼고
아직 형식화하지 않았다.
예정된 Lean 이름은 \texttt{\textbackslash lean} 대신 ``예정 Lean 이름''으로 적는다.

\section{기본 도구}

\begin{lemma}[꼭짓점 집합 넓히기]
\label{lem:map_iff}
\uses{def:IsPlanar}
\lean{Kuratowski.isPlanar_map_iff, Kuratowski.dartEquivMap, Kuratowski.RotationSystem.pushDart,
  Kuratowski.RotationSystem.pullDart, Kuratowski.ncard_isolated_map_add,
  Kuratowski.card_connectedComponent_map_add, Kuratowski.reachable_map_apply}
\leanok
\(W\)와 \(V\)가 유한하고 \(f : W \hookrightarrow V\)가 단사이면, \(W\) 위의 그래프 \(K\)에 대해
\(K.\mathrm{map}\,f\)가 평면인 것과 \(K\)가 평면인 것은 같다. \(K\)에 고립점이 있어도 된다.
\end{lemma}

\begin{proof}
\uses{def:rotation, def:eulerChar}
\leanok
\begin{itemize}
  \item 다트: \(K\)의 다트 \((a, b)\)와 \(K.\mathrm{map}\,f\)의 다트 \((f a, f b)\)는 일대일 대응한다. 이 대응은
        시작점을 \(f\)로 옮기고 다트 뒤집기와 교환한다.
  \item 회전 체계: 이 대응을 따라 회전 체계를 \(K\)에서 \(K.\mathrm{map}\,f\)로, 또 거꾸로 옮긴다. 면 순열은
        켤레가 되므로 순환 수 \(O\)가 같다 (\texttt{orbitCount\_permCongr}). 변 수 \(E\)도 같다.
  \item 수 세기: \(n = |V|\), \(m = |W|\)라 하자. \(f\)의 상에 없는 \(n - m\)개의 점은 \(K.\mathrm{map}\,f\)의
        고립점이고 각각 혼자 한 성분이다. 상 안에서는 \(f a\)와 \(f b\)가 이어진 것과 \(a\)와 \(b\)가
        이어진 것이 같다. 이것을 비교할 때는 점 \(a\)가 있으므로 \(W\)가 비어 있지 않고, 그러면 \(f\)의 왼쪽
        역함수가 \(K.\mathrm{map}\,f\)에서 \(K\)로 가는 그래프 준동형이 된다. \(W\)가 비어 있으면 그런 두 점이
        없고, \(V\)의 점은 모두 고립점이다. 따라서 \(I' + m = I + n\), \(c' + m = c + n\)이다.
        \(W\)나 \(V\)가 비어 있어도 된다.
  \item 그러면 \(\chi' - 2c' = (n - E + O + I') - 2c' = (m - E + O + I) - 2c = \chi - 2c\)이므로
        \(\chi' = 2c'\)와 \(\chi = 2c\)는 같은 조건이다.
\end{itemize}
\end{proof}

새 꼭짓점이 필요할 때 \(V\)를 \(\mathrm{Option}\,V\)로 넓히는 데 쓴다 (결정~D1, \ref{chap:plan}장).

\begin{lemma}[지지 집합이 4개 이하면 평면]
\label{lem:small}
\uses{def:IsPlanar}
\lean{Kuratowski.isPlanar_of_ncard_support_le_four}
\leanok
유한한 꼭짓점 타입 \(V\) 위의 그래프 \(G\)에서 지지 집합(고립점이 아닌 꼭짓점의 집합)의 크기가
4 이하면 \(G\)는 평면이다. 고립점 수에는 제한이 없고, \(V\)나 지지 집합이 비어 있어도 된다.
Lean의 \(G.\mathrm{support}\)가 이 집합이라는 것은 Mathlib의
\texttt{mem\_support\_iff\_not\_isIsolated}로 확인된다.
\end{lemma}

\begin{proof}
\uses{thm:K4, thm:mono, lem:map_iff}
\leanok
지지 집합 \(S = S_G\)에 대해 \(G_S\)를 \(G\)를 \(S\)로 제한한 그래프(꼭짓점 타입 \(S\))라 하자. 지지 집합 밖의
점은 고립점이므로 \(G = G_S.\mathrm{map}\,\iota\)이다 (\(\iota : S \hookrightarrow V\)는 포함 사상).
\(|S| \le 4\)이므로 단사 \(g : S \hookrightarrow \mathrm{Fin}\,4\)가 있고, \(\mathrm{Fin}\,4\) 위의 단순 그래프는 모두
완전 그래프 \(K_4\)의 부분 그래프이므로 \(G_S.\mathrm{map}\,g \le K_4\)이다.
\begin{enumerate}
  \item \(K_4\)는 평면이므로 (정리~\ref{thm:K4}) 부분 그래프 \(G_S.\mathrm{map}\,g\)도 평면이다
        (정리~\ref{thm:mono}).
  \item 보조정리~\ref{lem:map_iff}를 \(g\)에 쓰면 \(G_S\)가 평면이다.
  \item 보조정리~\ref{lem:map_iff}를 \(\iota\)에 쓰면 \(G = G_S.\mathrm{map}\,\iota\)가 평면이다.
\end{enumerate}
map의 방향에 주의한다. 사상은 \(S\)에서 \(\mathrm{Fin}\,4\)와 \(V\)로 가므로, \(G\)를 직접
\(K_4.\mathrm{map}\) 아래에 넣는 것이 아니라 \(G_S\)를 거쳐 두 번 옮긴다.
\end{proof}

\section{세분의 경로 표현}

\begin{definition}[경로 체계]
\label{def:paths}
\lean{Kuratowski.SubdivisionPaths, Kuratowski.walkInterior,
  Kuratowski.SubdivisionPaths.ofMapLE, Kuratowski.PathsExamples.k2InPath3}
\leanok
\(K\)(꼭짓점 집합 \(W\))의 \(G\) 안 \emph{경로 체계}는 다음으로 이루어진다.
\begin{itemize}
  \item 분기점 \(f : W \hookrightarrow V\)
  \item \(K\)의 변 \(ab\)마다 \(G\)의 경로 \(P_{ab}\) (\(f(a)\)에서 \(f(b)\)까지, \texttt{IsPath},
        \(P_{ba} = P_{ab}^{-1}\))
  \item 경로의 내부 꼭짓점은 \(f\)의 상 밖에 있다. 내부 집합은 경로의 \texttt{support}에서 두 끝점을
        뺀 것으로 정의한다 (\texttt{walkInterior}). \texttt{IsPath}이므로 끝점이 중간에 다시 나오지 않아,
        이 집합이 정확히 중간 꼭짓점들이다.
  \item 서로 다른 두 \emph{무방향} 변의 경로는 내부 꼭짓점을 공유하지 않는다.
        즉 \(\mathrm{Sym2}(a,b) \ne \mathrm{Sym2}(c,d)\)일 때 내부 집합이 서로소다.
        같은 변을 뒤집은 \(ab\), \(ba\)에는 이 조건을 적용하지 않는다.
\end{itemize}
경로는 Mathlib의 \texttt{G.Walk (f a) (f b)}로 저장하고, 단순성과 역방향 일치, 내부 회피와 서로소성을
각각 구조체 필드로 요구한다. 원래 그래프 \(K\)의 고립점도 분기점 단사에 포함된다.
이 정의 자체에는 유한성 가정이 필요 없다. 아래 동치 정리는 아직 형식화하지 않았다.
\texttt{ofMapLE}는 \(K.\mathrm{map}\,f \le G\)이면 각 변을 한 변짜리 경로로 하는 체계를 만든다.
\texttt{PathsExamples.lean}에서는 \(K_2\)의 한 변을 \(0,1,2\) 경로로 실현해 내부 \(\{1\}\)이 비지
않아도 허용됨을 확인한다. 또 유효한 단순 경로라도 내부에 분기점이 들어가거나 서로 다른 변의 후보 경로가 내부를 공유하면
해당 조건을 어김을 확인한다. 별 안의 경로 체계 전체에 대한 비존재 정리를 증명한 것은 아니다.
\end{definition}

\begin{theorem}[두 표현은 같다]
\label{thm:paths_iff}
\uses{def:subdiv, def:paths}
\planned{containsSubdivision\_iff\_paths}
\(G\)가 \(K\)의 세분을 포함하는 것은 \(K\)의 \(G\) 안 경로 체계가 있는 것과 같다.
\end{theorem}

\begin{proof}
(\(\Rightarrow\)) \texttt{IsSubdivision}에 대한 귀납법. 변 \(xy\)를 \(w\)로 나누면 \(xy\)가 들어 있던 경로
하나에 \(w\)가 끼어든다.

(\(\Leftarrow\)) 경로 내부 꼭짓점의 총수에 대한 귀납법. 내부 꼭짓점을 하나씩 \texttt{subdivideEdge}로
넣는다. 넣는 순간 그 꼭짓점은 아직 고립점이다.

이 동치는 아직 형식화하지 않았다. 귀납적인 세분 정의(정의~\ref{def:subdiv})가 이 경로 표현, 즉 표준적인
위상 소수 정의와 같다는 것은 쉬운 방향의 두 번째 리뷰에서 수학적으로 따로 확인했다.
\end{proof}

\begin{lemma}[공통 이웃 셋이면 \(TK_5\)]
\label{lem:TK5_of_common}
\uses{def:subdiv, def:K5K33, def:cycseq}
\planned{containsSubdivision\_K5\_of\_three\_common}
\(x \sim y\)이고, \(x\), \(y\)를 지나지 않는 사이클 \(C\)에 \(x\)와 \(y\)의 공통 이웃 \(z_1, z_2, z_3\)가
있으면 \(G\)는 \(K_5\)의 세분을 포함한다.
\end{lemma}

\begin{proof}
\uses{thm:paths_iff}
분기점은 \(x, y, z_1, z_2, z_3\)이다. 경로는 변 \(xy\), \(xz_i\), \(yz_i\)와 \(C\)의 세 호
\(z_i \cdots z_j\)다.
\end{proof}

\begin{lemma}[번갈아 놓이면 \(TK_{3,3}\)]
\label{lem:TK33_of_alt}
\uses{def:subdiv, def:K5K33, def:cycseq}
\planned{containsSubdivision\_K33\_of\_alternating}
\(x \sim y\)이고, 사이클 \(C\)(\(x\), \(y\)를 지나지 않음) 위에 \(u, x', w, x''\)가 이 순환 순서로 있다고
하자. \(u, w \in N(y)\)이고 \(x', x'' \in N(x)\)이면 \(G\)는 \(K_{3,3}\)의 세분을 포함한다.
\end{lemma}

\begin{proof}
\uses{thm:paths_iff}
두 편은 \(\{x, u, w\}\)와 \(\{y, x', x''\}\)이다. 경로는 변 5개(\(xy\), \(xx'\), \(xx''\), \(yu\), \(yw\))와
\(C\)의 네 호다.
\end{proof}
